\subsection{Additional geographical evidence and interpretation of disagreements}
\label{sec:litfollowup}
The initial 29 unmatched crop--country pairs are gaps in the recorded search, not proof that relevant research does
not exist. A targeted follow-up on 18 November 2026 identified regional future-yield evidence for six pairs and
regional future-suitability evidence for four more. These partial matches supplement the earlier assessment;
they are not ten independent validations of national yield predictions. Nineteen pairs still lack a verified
future comparison in this follow-up. Historical experiments and mechanism studies are recorded separately.
The original 39-row generated comparison remains the earlier assessment; its count is not silently replaced.

\begin{longtable}{L{0.22}L{0.36}L{0.42}}
\caption{Additional literature: directions and mechanisms, not numerical magnitudes.}\label{tab:litfollowup}\\
\toprule Pair & Evidence & Interpretation \\\midrule\endfirsthead
\toprule Pair & Evidence & Interpretation \\\midrule\endhead
Wheat--Egypt & Ali et al., Frontiers of Earth Science (2020), \href{https://doi.org/10.1007/s11707-019-0806-4}{doi:10.1007/s11707-019-0806-4}; North Nile Delta yield declines & Our historical-reference change is slightly negative but present-relative change slightly positive. Regional climate, cultivar, phenology and scenario differences matter. \\
Rice--Italy & Bocchiola, Agricultural Systems (2015), \href{https://doi.org/10.1016/j.agsy.2015.07.009}{doi:10.1016/j.agsy.2015.07.009}; mid-century CO$_2$-inclusive gains possible & Our present-relative direction is positive while the historical-reference direction is negative. This comparison is baseline-sensitive. \\
Rice--USA & Wilson et al., Global Change Biology (2026), \href{https://doi.org/10.1111/gcb.70679}{doi:10.1111/gcb.70679}; Texas observations/extrapolation & Combined genetic, climate and management effects differ from climate-only losses. Milling quality is distinct from micronutrient concentration. \\
Maize--South Africa & Ferreira et al., International Journal of Biometeorology (2023), \href{https://doi.org/10.1007/s00484-022-02392-1}{doi:10.1007/s00484-022-02392-1}; Limpopo losses & Our national result gives gains in both references. Regional drought, dynamic phenology and multiple climate models offer plausible explanations, not quantified attributions. \\
Soybean--South Africa & Olabanji et al., Climate (2021), \href{https://doi.org/10.3390/cli9010006}{doi:10.3390/cli9010006}; Olifants catchment losses & Our national result gives gains. Rainfall, soil-water accounting, geography and tested management differ. \\
Groundnut--USA & Zhen et al., Computers and Electronics in Agriculture (2024), \href{https://doi.org/10.1016/j.compag.2024.109105}{doi:10.1016/j.compag.2024.109105}; heterogeneous regional outcomes & Northern gains coexist with central/southern risks. A national mean hides this heterogeneity; cultivar traits also differ. \\
Rapeseed--China, USA, Brazil, South Africa & Borges et al., Scientific Reports (2023), \href{https://doi.org/10.1038/s41598-023-38910-3}{doi:10.1038/s41598-023-38910-3}; suitability losses in some regions & Suitability is not yield. Negative directions are qualitatively compatible in several cases; South Africa changes sign with our baseline. Our unsupported heat slope still limits interpretation. \\
\bottomrule
\end{longtable}

\paragraph{Reading limits.} Several additions are based on publisher abstracts or indexed results rather than full
methods. The downloaded Borges article was inspected for its suitability results. A seemingly Nigerian rice study
was rejected because its field sites are in Cote d'Ivoire. Wild-chickpea flowering forecasts, historical lentil
simulations, theses and market outlooks do not close cultivated-crop future-yield gaps. Source locators, reading
depth and all 29 pair decisions are recorded in \file{docs/evidence/LITERATURE\_FOLLOWUP.md} and
\file{docs/evidence/literature\_followup\_pairs.csv}.

\paragraph{Why conclusions differ.} First distinguish index, yield, production, suitability, nutrient supply and
intake. Then check baseline, period, geography and scenario. CO$_2$ fertilisation, season length, daily extremes,
soil storage, irrigation constraints and adaptation are documented model differences, but their individual
contributions have not been isolated. They should be described as plausible explanations, not established causes.

\paragraph{Markets and trade.} The project already uses static production, import/export and domestic-supply
accounting. It does not model endogenous future prices, trade, stocks, land allocation or household substitution.
These processes can change national availability and economic outcomes even when biological yield responses have
the same sign. Snyder et al., PLOS ONE (2020),
\href{https://doi.org/10.1371/journal.pone.0237918}{doi:10.1371/journal.pone.0237918}, demonstrate the importance of
international market effects for US agricultural financial outcomes. This supports the pathway, not attribution of
our crop-yield disagreements to imports. Quantitative attribution requires common-input model comparisons and
trade, stock, price and demand-response evidence that is not implemented here.
